\documentclass[12pt]{article}
\usepackage[utf8]{inputenc}
\usepackage[T1]{fontenc}
\usepackage[margin=1in]{geometry}
\usepackage{setspace}
\usepackage{xcolor}
\usepackage[colorlinks=true, urlcolor=blue!60!black, linkcolor=blue!60!black]{hyperref}
\setlength{\parindent}{0em}
\setlength{\parskip}{0.6em}
\renewcommand{\baselinestretch}{1.0}
\begin{document}
\begin{center}
{\LARGE \textit{Puerto Rico v. Branstad}, 483 U.S. 219 (1987)\par}
\vspace{0.5em}
{\large Case Brief --- Two-Page Edition\par}
\end{center}

\section{Memory Jogger}
Overruling \textit{Kentucky v.\ Dennison}, the Court unanimously held that federal courts may compel a governor to perform the mandatory duty to extradite a fugitive, here to Puerto Rico under the Extradition Act (\href{https://supreme.justia.com/cases/federal/us/483/219/}{483 U.S. 219}).

\section{Detailed Case Facts}
Ronald Calder, an Iowa native, was charged with first-degree murder in Puerto Rico after allegedly running over a pregnant woman with his car; he fled to Iowa. At the extradition hearing his lawyer claimed ``a white American man . . . could not receive a fair trial'' there, and Iowa Governors Ray and Branstad refused to surrender him.

\section{Procedural History}
Puerto Rico sought mandamus in the Southern District of Iowa; the district court dismissed under \textit{Dennison}, the Eighth Circuit ``[r]eluctantly'' affirmed, and the Supreme Court reversed.

\section{Judicial Votes}
Reversed, 9--0 (6--3 on Part II-B). Marshall wrote for the Court; O'Connor, with Powell, and Scalia concurred in part and in the judgment.

\section{Holding}
Federal courts may issue mandamus compelling an asylum governor to perform the ministerial duty of extradition; the Extradition Act, 18 U.S.C. \S~3182, gives Puerto Rico that right as a Territory.

\section{Analysis of Opinions}
Marshall placed \textit{Dennison} on the eve of the Civil War, when federal power was ``at its lowest ebb,'' and noted that Taney had found the duty ``merely ministerial'' yet unenforceable. Reaffirming that the Clause is mandatory, Marshall rejected Taney's coequal-sovereigns premise: federal courts routinely enforce federal duties against state officers, and a duty imposed by the Constitution itself requires no Tenth Amendment balancing. Section 3182 reaches any ``Territory,'' including Puerto Rico. O'Connor would have rested solely on the statute; Scalia joined only the statutory holding and noted that no party had questioned Congress's power to require extradition to a Territory.

\section{Examples: Future Applications}
\textbf{Same-Side:} a State's shield law cannot block extradition of a doctor who committed the charged act in the demanding State and fled. \textbf{Opposite-Side:} a telehealth prescriber who never entered the demanding State never ``fled from Justice,'' so no duty arises. \textbf{Fence-Sitter:} Puerto Rico's statutory demand after \textit{Printz}'s anti-commandeering rule.

\section{Critique}
The Court used a constitutional duty to bypass the Tenth Amendment while resting Puerto Rico's remedy on a statute, inviting tension with \textit{Printz}. Mandatory rendition also leaves asylum governors no humanitarian safety valve, the dilemma behind \textit{Dennison}.

\section{Key Quotations}
\begin{quote}
``\emph{Kentucky v. Dennison} is the product of another time. The conception of the relation between the States and the Federal Government there announced is fundamentally incompatible with more than a century of constitutional development. Yet this decision has stood while the world of which it was a part has passed away. We conclude that it may stand no longer.''
\end{quote}
\noindent --- Justice Marshall, Opinion of the Court, 483 U.S. at 230.

\end{document}
